\documentclass[10pt,a4paper,landscape]{article}
\usepackage[margin=0.45cm]{geometry}
\usepackage{amsmath,amssymb}
\usepackage{multicol}
\usepackage{enumitem}
\usepackage[T1]{fontenc}
\pagestyle{empty}
\setlength{\parindent}{0pt}
\setlength{\parskip}{1pt}
\setlength{\columnsep}{0.38cm}
\setlist[itemize]{leftmargin=*,nosep,topsep=1pt}
\setlist[enumerate]{leftmargin=*,nosep,topsep=1pt}
\renewcommand{\baselinestretch}{0.93}
\newcommand{\topic}[1]{\vspace{2pt}\textbf{\large #1}\par\hrule\vspace{2pt}}
\newcommand{\subtopic}[1]{\textbf{#1.}}
\newcommand{\R}{\mathbb R}
\newcommand{\Hcal}{\mathcal H}

\begin{document}
\fontsize{8.45}{9.15}\selectfont

\begin{center}
{\Large \textbf{ARIN7013 Exam Cheat Sheet, Page 1: CNN, Kernels, Linear Systems}}\\[-1mm]
\textbf{Use this as a hand-copy plan: formula first, then method steps, then final conclusion.}
\end{center}

\begin{multicols*}{4}

\topic{CNN And Convolution}
\subtopic{Cross-correlation}
For image $I$ and filter $K$,
\[
(I*K)_{ij}=\sum_{a,b} I_{i+a,j+b}K_{ab}.
\]
In CNN slides, ``convolution'' usually means cross-correlation: no filter flip.

\subtopic{Output shape}
For input $H\times W\times C_{\rm in}$, filter $F_h\times F_w$, padding $P_h,P_w$, stride $S_h,S_w$,
\[
H_{\rm out}=\Big\lfloor {H+2P_h-F_h\over S_h}\Big\rfloor+1,\quad
W_{\rm out}=\Big\lfloor {W+2P_w-F_w\over S_w}\Big\rfloor+1.
\]
Output channels $=C_{\rm out}$, the number of filters.

\subtopic{Weights}
\[
\#\text{conv weights}=F_hF_wC_{\rm in}C_{\rm out}
\]
Ignore bias if question says so; otherwise add $C_{\rm out}$.
Pooling has no trainable weights. Flatten only reshapes.
Fully connected: input dimension $\times$ output dimension.

\subtopic{CNN advantages}
Write: local receptive fields, sparse connectivity, parameter sharing, fewer parameters, grid/image inductive bias, translation-related robustness after convolution/pooling.

\subtopic{Residual CNN}
Plain block learns $F(X)$. Residual block:
\[
Y=X+F(X),\qquad F(X)=K_2*\sigma(K_1*X).
\]
Need same shape for addition. If shape/channel differs, use projection, e.g. $1\times1$ conv. Advantage: easier optimization, learns correction, better gradient flow.

\topic{Kernels And RKHS}
\subtopic{Positive definite kernel}
$K:X\times X\to\R$ is p.d. if for any $x_i$ and $c_i$,
\[
\sum_{i,j}c_ic_jK(x_i,x_j)\ge0.
\]
Equivalent Gram matrix $[K(x_i,x_j)]$ is positive semidefinite.

\subtopic{Inner product proof}
If $K(x,x')=\langle \phi(x),\phi(x')\rangle_{\Hcal}$, then
\[
\sum_{i,j}c_ic_jK(x_i,x_j)
=\Big\|\sum_i c_i\phi(x_i)\Big\|_{\Hcal}^2\ge0.
\]

\subtopic{Closure rules}
If $K_1,K_2$ are p.d., then $aK_1+bK_2$ for $a,b\ge0$, $K_1K_2$, pointwise limits, and $\exp(K_1)$ are p.d.
Polynomial $(1+\langle x,x'\rangle)^m$ is p.d.

\subtopic{Gaussian kernel proof}
\[
e^{-\gamma\|x-x'\|^2}
=e^{-\gamma\|x\|^2}e^{2\gamma x^Tx'}e^{-\gamma\|x'\|^2}.
\]
Use
\[
e^{2\gamma x^Tx'}=\sum_{m=0}^{\infty}{(2\gamma)^m\over m!}(x^Tx')^m,
\]
nonnegative sum of p.d. kernels, then multiply by $g(x)g(x')$.

\subtopic{Aronszajn theorem}
$K$ is p.d. iff there exists a Hilbert space $\Hcal$ and feature map $\phi$ such that
\[
K(x,x')=\langle \phi(x),\phi(x')\rangle_{\Hcal}.
\]
Exam sentence: kernels allow inner products in feature space without explicitly constructing features.

\subtopic{RKHS center distance}
For $\bar\phi={1\over n}\sum_j\phi(x_j)$,
\[
\|\phi(x_i)-\bar\phi\|_{\Hcal}^2
=K_{ii}-{2\over n}\sum_jK_{ij}
{1\over n^2}\sum_{j,l}K_{jl}.
\]
Average distance: average square root. Average squared distance: average this quantity.

\topic{KRR}
Objective in slides:
\[
\hat f=\arg\min_{f\in\Hcal}{1\over n}\sum_i |f(x_i)-y_i|^2+\beta\|f\|_{\Hcal}^2.
\]
Representer theorem:
\[
\hat f(x)=\sum_i \varsigma_iK(x_i,x).
\]
Let Gram $K_{ij}=K(x_i,x_j)$, $y=(y_i)$. Then fitted values $K\varsigma$ and $\|\hat f\|_{\Hcal}^2=\varsigma^TK\varsigma$.
\[
\varsigma=(K+\beta n I)^{-1}y
\]
up to null-space ambiguity; this choice is the standard answer. Final predictor must be written as the kernel expansion, not just the coefficient vector.

\topic{Linear Systems}
\subtopic{Gaussian elimination / LU}
Direct method; cost about $O(n^3)$, solve triangular systems in $O(n^2)$.
LU solve template:
\[
A=LU,\quad Ly=b,\quad Ux=y.
\]
Never write $x=A^{-1}b$ as a computational method unless asked conceptually.

\subtopic{Doolittle LU}
$A=LU$, $L_{kk}=1$.
\[
u_{kj}=a_{kj}-\sum_{s=1}^{k-1}\ell_{ks}u_{sj},\quad j=k,\ldots,n,
\]
\[
\ell_{ik}={a_{ik}-\sum_{s=1}^{k-1}\ell_{is}u_{sk}\over u_{kk}},\quad i=k+1,\ldots,n.
\]
Need nonzero pivot $u_{kk}$. Pivoting swaps rows when pivot is zero/small; partial pivoting uses largest available pivot in column.

\subtopic{2025 tridiagonal LU pattern}
For $A_n$ with diagonal $2$ and off-diagonal $-1$:
\[
U_{ii}={i+1\over i},\quad U_{i,i+1}=-1,\quad
L_{ii}=1,\quad L_{i+1,i}=-{i\over i+1}.
\]
Derive by small cases $n=3,4,5$:
$d_1=2$, $\ell_{i+1,i}=-1/d_i$, $d_{i+1}=2-1/d_i$.

\subtopic{Condition number}
\[
\kappa(A)=\|A\|\,\|A^{-1}\|.
\]
Large $\kappa(A)$ means small perturbations in $b$ or $A$ can cause large solution error. Ill-conditioned $\approx$ nearly singular.

\topic{Jacobi / Gauss-Seidel / SVD}
\subtopic{Jacobi}
Write $A=D+R$:
\[
x^{(k+1)}=D^{-1}(b-Rx^{(k)}),\quad
x_i^{(k+1)}={1\over a_{ii}}\Big(b_i-\sum_{j\ne i}a_{ij}x_j^{(k)}\Big).
\]
All RHS values are old values. Easy to parallelize.
Iteration matrix:
\[
G_J=-D^{-1}R=I-D^{-1}A,\quad e^{(k+1)}=G_Je^{(k)}.
\]
Converges for all $x^{(0)}$ and $b$ iff $\rho(G_J)<1$. Sufficient: strict diagonal dominance.

\subtopic{Gauss-Seidel}
Uses newest values immediately:
\[
x_i^{(k+1)}={1\over a_{ii}}\Big(b_i-\sum_{j<i}a_{ij}x_j^{(k+1)}
-\sum_{j>i}a_{ij}x_j^{(k)}\Big).
\]
Usually faster if convergent, less parallel.

\subtopic{Matrix multiplication}
$(n\times n)(n\times n)$: $O(n^3)$. $(m\times n)(n\times r)$: $O(mnr)$.

\subtopic{SVD}
\[
A=U\Sigma V^T,\quad A^TA v_i=\sigma_i^2 v_i,\quad u_i={Av_i\over \sigma_i}.
\]
Rank-$k$ approximation:
\[
A_k=\sum_{i=1}^k\sigma_i u_iv_i^T.
\]
Best in Frobenius/spectral norm:
\[
\|A-A_k\|_F^2=\sum_{i>k}\sigma_i^2,\qquad
\|A-A_k\|_2=\sigma_{k+1}.
\]

\end{multicols*}
\newpage

\begin{center}
{\Large \textbf{ARIN7013 Exam Cheat Sheet, Page 2: Eigen, PCA, ODE, ResNet, Neural ODE}}\\[-1mm]
\textbf{Most marks come from executing the algorithm with the lecture notation.}
\end{center}

\begin{multicols*}{4}

\topic{Eigenvalue Algorithms}
\subtopic{Rayleigh quotient}
For symmetric $A$,
\[
R_A(v)={v^TAv\over v^Tv},\quad \text{if }\|v\|=1,\ R_A(v)=v^TAv.
\]
Eigenvectors are stationary points. Largest eigenvalue maximizes $R_A(v)$ over unit vectors.

\subtopic{Power iteration}
For dominant eigenvalue $|\lambda_1|>|\lambda_2|$ and $x_0=a_1u_1+\cdots$, $a_1\ne0$:
\[
z^{(k)}=Av^{(k-1)},\quad v^{(k)}={z^{(k)}\over \|z^{(k)}\|},\quad
\omega^{(k)}=(v^{(k)})^TAv^{(k)}.
\]
Proof template:
\[
A^kx_0=a_1\lambda_1^k\left(u_1+\sum_{i\ge2}{a_i\over a_1}\left({\lambda_i\over\lambda_1}\right)^ku_i\right).
\]
Convergence factor roughly $|\lambda_2/\lambda_1|$.

\subtopic{Inverse iteration}
Choose shift $\nu$ near desired eigenvalue:
\[
(A-\nu I)z^{(k)}=v^{(k-1)},\quad
v^{(k)}={z^{(k)}\over\|z^{(k)}\|}.
\]
Targets eigenvalue closest to $\nu$ because eigenvalues of $(A-\nu I)^{-1}$ are $1/(\lambda_i-\nu)$.

\subtopic{Rayleigh quotient iteration}
\[
\omega^{(k-1)}=(v^{(k-1)})^TAv^{(k-1)},
\]
\[
(A-\omega^{(k-1)}I)z^{(k)}=v^{(k-1)},\quad
v^{(k)}={z^{(k)}\over\|z^{(k)}\|}.
\]
For symmetric $A$, sufficiently close to simple eigenpair: cubic convergence.

\topic{Ordinary PCA}
\subtopic{Core notation}
Data $x_i\in\R^d$. Center unless question says not to:
\[
\tilde x_i=x_i-m,\quad m={1\over n}\sum_i x_i.
\]
Covariance:
\[
V={1\over n}\sum_i \tilde x_i\tilde x_i^T={1\over n}X^TX
\]
if rows are samples. If columns are samples, use $V={1\over n}XX^T$.

\subtopic{Loading vs score}
Principal component asked in exams usually means loading vector:
\[
Vh_j=\omega_jh_j,\quad \|h_j\|=1.
\]
PC score for datum $x$:
\[
y_j=\langle x,h_j\rangle \quad \text{or}\quad y_j=\langle \tilde x,h_j\rangle.
\]
Do not answer with score when asked for loading/component.

\subtopic{Projection and MSE}
Projection onto first $s$ PCs:
\[
P_s(x_i)=\sum_{j=1}^s\langle x_i,h_j\rangle h_j.
\]
Average reconstruction MSE:
\[
{1\over n}\sum_i\left\|x_i-P_s(x_i)\right\|^2
=\sum_{j>s}\omega_j.
\]
Information preserved:
\[
{\sum_{j=1}^s\omega_j\over \sum_{j=1}^d\omega_j}\times100\%.
\]

\subtopic{If $d$ is large}
Use $XX^T/n$ or $X^TX/n$ trick. Nonzero eigenvalues match up to dimension convention. For SVD $X=U\Sigma V^T$, PCA directions are singular vectors; variances are $\sigma_i^2/n$.

\topic{Kernel PCA}
\subtopic{Setup}
Feature map $\Psi:X\to\Hcal$, kernel
\[
K(x,x')=\langle \Psi(x),\Psi(x')\rangle_{\Hcal},\quad
K_{ij}=K(x_i,x_j).
\]
Assume centered in feature space unless question says no centering. If no centering, use raw Gram matrix.

\subtopic{Eigenproblem}
Slide covariance:
\[
V={1\over n}\sum_i \Psi(x_i)\Psi(x_i)^T,\quad
v_j=\sum_i a_{ji}\Psi(x_i).
\]
Here $v_j\in\Hcal$ is the loading/eigenfunction. It is not an $n$-vector.
Then
\[
Ka_j=n\omega_j a_j,\qquad a_j^TKa_j=1.
\]
Exam calculation route:
\[
Kq_j=\rho_jq_j,\quad \|q_j\|=1,\quad
a_j={q_j\over\sqrt{\rho_j}},\quad \rho_j=n\omega_j.
\]
Then actual RKHS PC:
\[
v_j=\sum_i a_{ji}\Psi(x_i).
\]
$q_j,a_j\in\R^n$ are coefficient vectors only.

\subtopic{Projection and MSE}
For new $x$:
\[
y_j(x)=\langle \Psi(x),v_j\rangle_{\Hcal}
=\sum_i a_{ji}K(x_i,x).
\]
Training score vector $=Ka_j$, not the loading.
Using first $s$ kernel PCs:
\[
\text{MSE}={1\over n}\left(\operatorname{tr}(K)-\sum_{j=1}^s\rho_j\right)
=\sum_{j>s}\omega_j.
\]
Common trap: answer must be $v_j\in\Hcal$, represented by $a_j$, not score vector $Ka_j$.

\topic{ODE Solvers}
\subtopic{IVP}
\[
x'(t)=f(t,x),\quad x(t_0)=x_0,\quad t_i=t_0+ih.
\]

\subtopic{Theta method}
\[
x_{i+1}=x_i+h[(1-\theta)f(t_i,x_i)+\theta f(t_{i+1},x_{i+1})].
\]
$\theta=0$: Forward Euler, explicit.
\[
x_{i+1}=x_i+hf(t_i,x_i).
\]
$\theta=1$: Backward Euler, implicit.
\[
x_{i+1}=x_i+hf(t_{i+1},x_{i+1}).
\]
$\theta=1/2$: trapezoidal, generally implicit,
\[
x_{i+1}=x_i+{h\over2}[f(t_i,x_i)+f(t_{i+1},x_{i+1})].
\]

\subtopic{2024 explicit improved Euler scheme}
\[
x_i=x_{i-1}+{h\over2}\{f(t_{i-1},x_{i-1})
 +f(t_i,x_{i-1}+hf(t_{i-1},x_{i-1}))\}.
\]
It is explicit: RHS uses known $x_{i-1}$ only.

\subtopic{Error estimate template}
Define exact/numerical error
\[
e_i=x(t_i)-x_i.
\]
Local truncation/consistency error, example theta $1/2$:
\[
T_i={x(t_{i+1})-x(t_i)\over h}
-{1\over2}[f(t_i,x(t_i))+f(t_{i+1},x(t_{i+1}))].
\]
Taylor with bounded derivatives gives $|T_i|\le Ch^2$ for trapezoidal/improved Euler, and $|T_i|\le Ch$ for Forward Euler.

Subtract numerical scheme from exact consistency relation. Use Lipschitz:
\[
|f(t,x)-f(t,y)|\le L|x-y|.
\]
Then get recurrence like
\[
|e_{i+1}|\le (1+Ch)|e_i|+Ch^{p+1}.
\]
Discrete Gronwall:
\[
\max_i |e_i|\le C h^p.
\]
Forward Euler global order $1$; trapezoidal/improved Euler global order $2$.

\topic{ResNet And Neural ODE}
\subtopic{ResNet as Euler}
Forward Euler for hidden state:
\[
u_{l+1}=u_l+h f(u_l,\theta_l).
\]
ResNet block:
\[
u_{l+1}=u_l+F(u_l,\theta_l).
\]
Interpret depth as time steps; residual connection improves gradient flow and stability.

\subtopic{Neural ODE}
\[
{dx(t)\over dt}=f(x(t),t,\theta),\quad
x(t_1)=\operatorname{ODESolve}(f,x(t_0),t_0,t_1).
\]
Continuous-depth model; solver chooses steps adaptively.

\subtopic{Adjoint gradient}
Loss $L=L(x(t_1))$. Define adjoint
\[
a(t)={\partial L\over \partial x(t)}.
\]
Terminal condition:
\[
a(t_1)={\partial L\over \partial x(t_1)}.
\]
Backward adjoint ODE:
\[
{da(t)\over dt}=-a(t)^T{\partial f(x(t),t,\theta)\over \partial x}.
\]
Parameter gradient:
\[
{\partial L\over\partial\theta}
=\int_{t_0}^{t_1}a(t)^T{\partial f(x(t),t,\theta)\over\partial\theta}\,dt.
\]
Memory comparison: ResNet/backprop stores many layer activations; Neural ODE adjoint can recompute trajectory backward, lower memory but possible numerical instability/recomputation cost.

\topic{Last-Minute Traps}
\begin{itemize}
\item PCA component/loading $\ne$ PC score. Loading is eigenvector.
\item Kernel PCA: answer component is $v\in\Hcal$; normalize $a^TKa=1$.
\item KRR: write final function $\hat f(x)=\sum_i\varsigma_iK(x_i,x)$.
\item Jacobi uses all old values; Gauss-Seidel uses new values immediately.
\item CNN: shape formula before weight count; filter weights are shared.
\item Do not center data if exam explicitly says no centering.
\item ODE error proof: truncation error $\to$ recurrence $\to$ Gronwall/order.
\item Residual addition requires matching dimensions.
\end{itemize}

\end{multicols*}
\end{document}
